%% The body of the paper — every section lives here.
%%
%% THIS FILE IS YOURS. `make set-template` rewrites main.tex but never touches
%% this file, so the same prose compiles under every template. Two rules keep
%% that true (see tools/templates/README.md):
%%
%%   * cite with \citep / \citet — every template loads natbib
%%   * anything you \begin{} must come from researchcommon.sty, not from one
%%     template's preamble

\section{Introduction}
\label{sec:introduction}

Experiment code and the paper written about it normally live in separate
repositories, and the numbers travel between them by hand. They then go stale
silently: nothing connects the figure in the PDF to the runs it summarizes, and
a table copied out of a notebook keeps its authority long after the sweep that
produced it was rerun. This repository closes that gap without merging the two
codebases: the paper fetches its numbers from the experiment tracker, records
which runs it fetched, and fails to build when the two disagree.

\begin{contributions}
\item A committed, human-readable cache of experiment-tracker results, with
  the exact run identifiers behind it pinned in a lockfile, so a paper's
  numbers are reproducible without access to the tracker.
  \label{c:cache}
\item Figures and tables generated from that cache with recorded provenance,
  which makes staleness a build failure rather than a reading error.
  \label{c:provenance}
\item A document directory that compiles standalone---on Overleaf, and as a
  flattened submission bundle---and whose prose survives a change of venue
  template.
  \label{c:portable}
\end{contributions}

\section{Method}
\label{sec:method}

% NOTE: `% NOTE:` comments are author instructions for prose that is planned but
% not yet written. Agents must treat the nearest NOTE as the specification for
% what belongs there, and must never delete or rewrite one. See ../../AGENTS.md.

% NOTE: describe the actual method here. This placeholder section exists only so
% the example compiles to a plausible short paper.

The synthetic study behind \cref{fig:learning_curve} stands in for a real
experiment. Deep reinforcement learning from pixels \citep{mnih2015human} is
sample-inefficient, which motivated episodic control methods
\citep{blundell2016modelfree,pritzel2017neural}. \citet{oquab2024dinov2} showed
that self-supervised features transfer broadly; replacing an episodic
controller's state encoder with such features is the kind of question this
repository layout is built to support.

\begin{figure}[!t]
  \centering
  \includegraphics[width=\linewidth]{architecture.pdf}
  \caption{Example schematic: a world-model control loop. The source is
    \texttt{figures/architecture.drawio}; regenerate the export with
  \texttt{make drawio}.}
  \label{fig:architecture}
\end{figure}

Hand-authored diagrams live beside the data-figure configs, as
\texttt{.drawio} sources. \Cref{fig:architecture} is one: it is drawn by hand,
not fetched from the tracker, and its PDF is a generated export.

\section{Results}
\label{sec:results}

Neither \cref{fig:learning_curve} nor \cref{tab:final_scores} is drawn or typed
by hand. Both are rendered from \texttt{data/learning\_curves.csv} and
\texttt{data/final\_scores.csv}, which \texttt{make fetch} writes from the runs
pinned in \texttt{data/wandb.lock.yml}. Each output records the content hash of
the file it read, so changing the data without rebuilding fails
\texttt{make check-generated} instead of quietly shipping a stale result.

\begin{figure}[!t]
  \centering
  \includegraphics[width=\linewidth]{learning_curve.pdf}
  \caption{Learning curves for the synthetic example study: mean over five
    seeds, shaded band is one standard deviation. Regenerate with
  \texttt{make plots}.}
  \label{fig:learning_curve}
\end{figure}

\begin{table}[!t]
  \centering
  \caption{Final score after training, mean over five seeds. The tabular is
  generated; the caption, label and float placement are not.}
  \label{tab:final_scores}
  %% GENERATED FILE — do not edit.
%%
%% Written from tables/final_scores.yml
%% out of the committed W&B cache. Rebuild with: make tables
%%
%% Editing this file by hand puts the paper's numbers out of step with the
%% runs behind them, and `make check-generated` will not catch it — it checks
%% the inputs, not the output. Change the config instead.
%%
%% This is a bare tabular: wrap it in \begin{table} with your caption and
%% label in text.tex.

\begin{tabular}{lrrr}
\toprule
Method & Score & SD & Seeds \\
\midrule
pretrained & 0.86 & 0.01 & 5 \\
scratch & 0.67 & 0.01 & 5 \\
baseline & 0.62 & 0.01 & 5 \\
\bottomrule
\end{tabular}

\end{table}

\begin{figure}[!t]
  \centering
  \includegraphics[width=\linewidth]{dreamer_optimisations_ablation.pdf}
  \caption{DreamerV3 runtime-optimisation ablation on Atari-100k. Rendered from
    \texttt{data/figures/dreamer\_optimisations\_ablation.csv}, pinned to the
  runs in \texttt{data/wandb.lock.yml}. Regenerate with \texttt{make plots}.}
  \label{fig:dreamer_optimisations_ablation}
\end{figure}

\section{Conclusion}
\label{sec:conclusion}

Delete this paper and write your own. Keep the structure: one bibliography and
one cache as sources of truth, generated files committed where the compiler
needs them, and a check that tells you when the two have drifted apart.
